% Barvinok-Pataki for Order Unit Spaces
% Adversarial formalisation report
%
% To compile: pdflatex main.tex && bibtex main && pdflatex main.tex && pdflatex main.tex

\documentclass[11pt,a4paper]{article}

%% ============================================================
%% Packages
%% ============================================================
\usepackage[utf8]{inputenc}
\usepackage[T1]{fontenc}
\usepackage{lmodern}
\usepackage[margin=1in]{geometry}
\usepackage{amsmath,amssymb,amsthm}
\usepackage{mathtools}
\usepackage{mathrsfs}
\usepackage{enumitem}
\usepackage{booktabs}
\usepackage{longtable}
\usepackage{tikz}
\usetikzlibrary{trees,arrows.meta}
\usepackage{hyperref}
\hypersetup{
    colorlinks=true,
    linkcolor=blue!70!black,
    citecolor=blue!70!black,
    urlcolor=blue!70!black
}
\usepackage[capitalise,nameinlink]{cleveref}

%% ============================================================
%% Theorem Environments
%% ============================================================
\theoremstyle{plain}
\newtheorem{theorem}{Theorem}[section]
\newtheorem{lemma}[theorem]{Lemma}
\newtheorem{proposition}[theorem]{Proposition}
\newtheorem{corollary}[theorem]{Corollary}

\theoremstyle{definition}
\newtheorem{definition}[theorem]{Definition}
\newtheorem{example}[theorem]{Example}

\theoremstyle{remark}
\newtheorem{remark}[theorem]{Remark}
\newtheorem*{remark*}{Remark}
\newtheorem*{caveat*}{Caveat}

%% ============================================================
%% Custom Commands
%% ============================================================
\newcommand{\N}{\mathbb{N}}
\newcommand{\Z}{\mathbb{Z}}
\newcommand{\R}{\mathbb{R}}
\newcommand{\C}{\mathbb{C}}
\newcommand{\Hilbert}{\mathcal{H}}
\newcommand{\BH}{\mathcal{B}(\Hilbert)}
\newcommand{\sa}{\mathrm{sa}}
\newcommand{\GNS}{\mathrm{GNS}}
\newcommand{\tr}{\mathrm{Tr}}
\newcommand{\rank}{\mathrm{rank}}
\newcommand{\aff}{\mathrm{aff}}
\newcommand{\inner}[2]{\langle #1, #2 \rangle}
\newcommand{\norm}[1]{\|#1\|}
\newcommand{\abs}[1]{|#1|}
\DeclareMathOperator{\dimR}{dim_{\R}}
\DeclareMathOperator{\Spec}{Spec}

%% ============================================================
%% Status box
%% ============================================================
\usepackage{tcolorbox}
\newtcolorbox{statusbox}[1][]{colback=yellow!10!white, colframe=orange!80!black,
  title={\textbf{Verification Status}}, fonttitle=\bfseries, #1}

%% ============================================================
%% Title
%% ============================================================
\title{\textbf{A Barvinok--Pataki Theorem for\\Order Unit Spaces}\\[0.5em]
       \large Adversarial Formalisation Report}
\author{Generated via Adversarial Proof Framework (\texttt{af})\\[0.3em]
        \small Prover--verifier protocol with opus-level adversarial verification}
\date{\today}

%% ============================================================
\begin{document}
\maketitle

%% ============================================================
%% Status Warning
%% ============================================================
\begin{statusbox}
This document reports on an \emph{adversarial formalisation} of the
Barvinok--Pataki theorem for order unit spaces, carried out using the
\texttt{af} (Adversarial Proof Framework) CLI tool.

\medskip\noindent
\textbf{Current state:} The proof tree contains 64 nodes. Of these,
\textbf{15 have been validated} by independent adversarial verifiers, and
\textbf{49 remain pending} (awaiting verifier acceptance). All 148 challenges
raised during verification have been resolved, but the resolved nodes have
\emph{not yet been re-verified}. The quality score is 69.4/100.

\medskip\noindent
\textbf{This is not a finished paper.} It is a structured proof skeleton
with adversarial annotations. Four genuine mathematical errors were found
and corrected during verification (see \cref{sec:errors}). The proof
architecture is complete, but full validation requires a second verifier pass.
\end{statusbox}

\tableofcontents
\newpage

%% ============================================================
%% 1. INTRODUCTION
%% ============================================================
\section{Introduction}
\label{sec:introduction}

This paper establishes an abstract Barvinok--Pataki theorem for state spaces
of $C^*$-algebras and Archimedean order unit spaces, unifying three previously
disconnected literatures:

\begin{enumerate}
  \item \textbf{SDP optimisation} --- the Barvinok--Pataki rank bounds
        for extreme points of spectrahedra
        \cite{barvinok2001,pataki1998,pataki2000};
  \item \textbf{Abstract convex geometry} --- the Weis--Shirokov face-dimension
        theorem for constrained convex sets \cite{weisshirokov2021};
  \item \textbf{Operator algebras} --- the Alfsen--Shultz face-commutant
        correspondence for state spaces \cite{alfsen2001,alfsen2003}.
\end{enumerate}

The key insight is that the face-commutant correspondence converts the
geometric face-dimension bound of Weis--Shirokov into an algebraic bound
on the commutant of the GNS representation, yielding a multiplicity bound
that simultaneously recovers all known special cases.


%% ============================================================
%% 2. MAIN RESULT
%% ============================================================
\section{Main Result}
\label{sec:main-result}

\begin{theorem}[Barvinok--Pataki for Order Unit Spaces]
\label{thm:main}
Let $\mathcal{A}$ be a $C^*$-algebra (or, more generally, an Archimedean
order unit space), and let $a_1, \ldots, a_m \in \mathcal{A}^{\sa}$ be
self-adjoint elements. Define the affinely constrained state space
\[
  C = \bigl\{\varphi \in S(\mathcal{A}) : \varphi(a_j) = \beta_j,\;
  j = 1, \ldots, m\bigr\}
\]
(or with inequality constraints $\varphi(a_j) \le \beta_j$, or mixed).
If $\varphi$ is an extreme point of $C$, then:
\begin{enumerate}[label=(\Roman*)]
  \item $\dimR\bigl(\pi_\varphi(\mathcal{A})'^{\,\sa}\bigr) \le m + 1$,
        where $\pi_\varphi(\mathcal{A})'$ is the commutant of the GNS
        representation.
  \item Equivalently, if
        $\pi_\varphi \cong \bigoplus_{j=1}^k \sigma_j^{\oplus n_j}$
        is the decomposition into pairwise inequivalent irreducible
        representations $\sigma_j$ with multiplicities $n_j$, then
        \[
          \sum_{j=1}^k n_j^2 \le m + 1.
        \]
\end{enumerate}
\end{theorem}

\begin{remark*}
The theorem extends to $*$-algebras equipped with an Archimedean quadratic
module $M$, in which case $S(\mathcal{A})$ is replaced by the
$M$-positive state space $S_M$ and compactness follows from Tychonoff's
theorem rather than Banach--Alaoglu.
\end{remark*}


%% ============================================================
%% 3. PROOF ARCHITECTURE
%% ============================================================
\section{Proof Architecture}
\label{sec:proof-architecture}

The proof proceeds in three steps, each corresponding to a major node
in the proof tree.

\subsection{Step 1: Face-Commutant Correspondence (Alfsen--Shultz)}
\label{sec:face-commutant}

\begin{proposition}[Node 1.1]
\label{prop:face-commutant}
Let $\mathcal{A}$ be a $C^*$-algebra and $\varphi \in S(\mathcal{A})$
a state with GNS triple $(\pi_\varphi, \Hilbert_\varphi, \xi_\varphi)$.
The face $F_{S(\mathcal{A})}(\varphi)$ generated by $\varphi$ in the
state space is affinely isomorphic to
\[
  D = \bigl\{T \in \pi_\varphi(\mathcal{A})'^{\,+} :
  \inner{T\xi_\varphi}{\xi_\varphi} = 1\bigr\}
\]
via the map $\Phi(T) = \omega_T$ where
$\omega_T(a) = \inner{\pi_\varphi(a) T \xi_\varphi}{\xi_\varphi}$.
Consequently,
\begin{equation}
\label{eq:face-dim}
  \dim\bigl(F_{S(\mathcal{A})}(\varphi)\bigr)
  = \dimR\bigl(\pi_\varphi(\mathcal{A})'^{\,\sa}\bigr) - 1.
\end{equation}
\end{proposition}

\begin{proof}[Proof sketch]
The proof establishes five properties of $\Phi$:
\begin{enumerate}[label=(\roman*)]
  \item \textbf{Well-definedness} (Node 1.1.1): $\omega_T$ is a state.
        Positivity uses $T \in \pi_\varphi(\mathcal{A})'$ to commute $T$
        past $\pi_\varphi(a)$:
        $\omega_T(a^*a) = \inner{T\pi_\varphi(a)\xi_\varphi}
        {\pi_\varphi(a)\xi_\varphi} \ge 0$.
  \item \textbf{Affinity} (Node 1.1.2): follows from linearity of the
        inner product (with real convex coefficients).
  \item \textbf{Image in face} (Node 1.1.3): $\omega_T \le \norm{T}\varphi$
        as positive functionals, so $\omega_T \in F_{S(\mathcal{A})}(\varphi)$
        by the face characterisation
        ($\psi \in F(\varphi) \Leftrightarrow \psi \le \lambda\varphi$
        for some $\lambda > 0$).
  \item \textbf{Injectivity} (Node 1.1.4): uses cyclicity of $\xi_\varphi$
        and the fact that $T$ commutes with $\pi_\varphi(\mathcal{A})$.
  \item \textbf{Surjectivity} (Node 1.1.5): by the Radon--Nikodym theorem
        for states (Alfsen--Shultz \cite{alfsen2001}, Prop.\ 5.31).
\end{enumerate}
For the dimension formula \eqref{eq:face-dim}, $D$ is the cross-section
$\pi_\varphi(\mathcal{A})'^{\,+} \cap f^{-1}(1)$ where
$f(T) = \inner{T\xi_\varphi}{\xi_\varphi}$. Since the positive cone
spans $\pi_\varphi(\mathcal{A})'^{\,\sa}$ and has nonempty interior
(containing $I$), the affine span of $D$ equals $f^{-1}(1)$, giving
$\dim(D) = \dimR(\pi_\varphi(\mathcal{A})'^{\,\sa}) - 1$
(Node 1.1.6).
\end{proof}

\subsection{Step 2: Abstract Face-Dimension Bound (Weis--Shirokov)}
\label{sec:weis-shirokov}

\begin{proposition}[Node 1.2]
\label{prop:weis-shirokov}
Let $K$ be a compact convex set in a locally convex space, and let
$f_1, \ldots, f_m : K \to \R \cup \{+\infty\}$ be lower semicontinuous
generalised affine maps. Define $C = \{x \in K : f_j(x) \le \beta_j\}$
(or $= \beta_j$, or mixed). For any $x \in C$:
\begin{equation}
\label{eq:ws-bound}
  \dim\bigl(F_K(x)\bigr) \le \dim\bigl(F_C(x)\bigr) + m.
\end{equation}
In particular, if $x$ is an extreme point of $C$, then
$\dim(F_K(x)) \le m$.
\end{proposition}

This is taken as an external result: \cite{weisshirokov2021}, Theorem~2
(single constraint, level sets), Theorem~4 (sublevel sets), and
Corollary~5 (multi-constraint induction). The multi-constraint case
follows by iterating the single-constraint bound: define
$K_0 = K$, $K_j = \{x \in K_{j-1} : f_j(x) \le \beta_j\}$,
and telescope $\dim(F_{K_{j-1}}(x)) \le \dim(F_{K_j}(x)) + 1$.

\begin{remark}
The naive ``constraint propagation'' argument --- if $y \in F_K(x)$ then
$f(y) = \beta$ --- does \emph{not} work for sublevel sets, since $f(z)$
can compensate for $f(y) > \beta$ in a convex combination.
Weis--Shirokov's proof handles this through a careful face-structure
analysis.
\end{remark}

\subsection{Step 3: Combining the Two Bounds}
\label{sec:combining}

\begin{proof}[Proof of \cref{thm:main}]
The argument proceeds in three lines (Nodes 1.3.1--1.3.4):
\begin{enumerate}
  \item \textbf{Hypothesis verification} (Node 1.3.1):
        For $C^*$-algebras, $K = S(\mathcal{A})$ is weak-$*$ compact
        (Banach--Alaoglu), and $f_j(\psi) = \psi(a_j)$ are continuous
        affine (by definition of the weak-$*$ topology).
        The face-commutant correspondence (\cref{prop:face-commutant})
        and the Weis--Shirokov bound (\cref{prop:weis-shirokov}) both apply.
  \item \textbf{Geometric bridge} (Node 1.3.2):
        By \eqref{eq:face-dim},
        $\dim(F_{S(\mathcal{A})}(\varphi))
        = \dimR(\pi_\varphi(\mathcal{A})'^{\,\sa}) - 1$.
  \item \textbf{Extremality bound} (Node 1.3.3):
        Since $\varphi$ is extreme in $C$,
        $\dim(F_C(\varphi)) = 0$, so by \eqref{eq:ws-bound},
        $\dim(F_{S(\mathcal{A})}(\varphi)) \le m$.
  \item \textbf{Conclusion} (Node 1.3.4):
        $\dimR(\pi_\varphi(\mathcal{A})'^{\,\sa}) - 1 \le m$,
        hence $\dimR(\pi_\varphi(\mathcal{A})'^{\,\sa}) \le m + 1$.
\end{enumerate}
For part (II), by the Artin--Wedderburn theorem (Node 1.4),
$\pi_\varphi(\mathcal{A})' \cong \bigoplus_{j=1}^k M_{n_j}(\C)$,
so $\dimR(\pi_\varphi(\mathcal{A})'^{\,\sa})
= \sum_{j=1}^k n_j^2 \le m + 1$.
\end{proof}


%% ============================================================
%% 4. RECOVERY OF CLASSICAL RESULTS
%% ============================================================
\section{Recovery of Classical Results}
\label{sec:recovery}

\subsection{Abelian Case: Carath\'eodory's Theorem}
\label{sec:caratheodory}

When $\mathcal{A} = C(X)$ for a compact Hausdorff space $X$, states
are probability measures $\mu$ on $X$. All irreducible representations
are one-dimensional (point evaluations $\mathrm{ev}_x$), and all
multiplicities equal~1 (since the commutant is abelian).
For a finitely-supported measure $\mu = \sum_{j=1}^k \lambda_j \delta_{x_j}$,
the GNS commutant is $\C^k$, giving $\sum n_j^2 = k$.

The main bound yields $k \le m + 1$: extreme points of the set of
probability measures satisfying $m$ moment constraints are discrete
measures supported on at most $m + 1$ points. This is the
\emph{Richter--Rogosinski theorem}, the measure-theoretic form of
Carath\'eodory's theorem.

\subsection{Matrix Algebra: Classical Barvinok--Pataki}
\label{sec:classical-bp}

When $\mathcal{A} = M_n(\C)$, states are density matrices
$\rho \in M_n(\C)$ with $\rho \ge 0$, $\tr(\rho) = 1$.
The unique irreducible representation is the identity on $\C^n$,
so a rank-$r$ state has GNS $\pi \cong \mathrm{id}^{\oplus r}$
with commutant $M_r(\C)$.
The bound gives $r^2 \le m + 1$.

The relationship to the standard SDP formulation $r^2 \le m$ (for the
Hermitian PSD cone) is explained by the trace normalisation: the state
space $S(M_n(\C))$ already incorporates $\tr(\rho) = 1$, which in
the standard formulation is counted among the $m$ constraints.

\begin{remark}
The full classification by real division algebras (Hurwitz) yields
three cases:
\begin{enumerate}[label=(\arabic*)]
  \item \emph{Real:} $\mathcal{A} = M_n(\R)$ with transpose involution,
        $\dim = r(r+1)/2$, bound $r(r+1)/2 \le m + 1$;
  \item \emph{Complex:} $\mathcal{A} = M_n(\C)$ with conjugate transpose,
        $\dim = r^2$, bound $r^2 \le m + 1$;
  \item \emph{Quaternionic:} $\mathcal{A} = M_n(\mathbb{H})$ with
        conjugate transpose, $\dim = r(2r-1)$, bound $r(2r-1) \le m + 1$.
\end{enumerate}
\end{remark}

\subsection{Block-Diagonal Algebras: Mixed Bound}
\label{sec:block-diagonal}

When $\mathcal{A} = \bigoplus_{i=1}^N M_{d_i}(\C)$, a state
$\varphi = \sum_i \lambda_i \tr(\rho_i \cdot)$ with active blocks
$I = \{i : \lambda_i > 0\}$ has GNS
$\pi_\varphi \cong \bigoplus_{i \in I} \mathrm{id}_{d_i}^{\oplus r_i}$
where $r_i = \rank(\rho_i)$. The bound becomes
\[
  \sum_{i \in I} r_i^2 \le m + 1.
\]
This simultaneously constrains ranks within blocks \emph{and} the number
of active blocks, interpolating between Carath\'eodory (all $d_i = 1$)
and Barvinok--Pataki ($N = 1$).

\begin{example}
For $\mathcal{A} = M_2(\C) \oplus M_2(\C)$ with $m = 4$ constraints:
the bound gives $r_1^2 + r_2^2 \le 5$. This rules out
$(r_1, r_2) = (2, 2)$ (which would require $\sum r_i^2 = 8 > 5$).
Applying classical BP to each block independently would allow full rank
in both blocks. The cross-block constraint is genuinely new.
\end{example}

\subsection{Infinite-Dimensional Case: Weis--Shirokov Recovery}
\label{sec:weis-shirokov-recovery}

When $\mathcal{A} = \BH$ for separable $\Hilbert$, normal states
correspond to density operators $\rho$ with $\tr(\rho) = 1$.
A rank-$r$ state has commutant $M_r(\C)$ and the bound gives
$r^2 \le m + 1$, i.e., $r \le \sqrt{m + 1}$.
This recovers the Weis--Shirokov bound \cite{weisshirokov2021}, Theorem~3.

For \emph{unbounded} Hamiltonians $H_j$ (self-adjoint operators affiliated
with $\BH$), the constraint $f(\rho) = \tr(H_j\rho)$ defines a lower
semicontinuous generalised affine map (as a supremum of continuous affine
functions via spectral truncation), and the Weis--Shirokov theorem
applies directly via Theorems~2 and~4 of \cite{weisshirokov2021}.


%% ============================================================
%% 5. FACE DIMENSION GAP AND AUTOMATIC PURITY
%% ============================================================
\section{Face Dimension Gap and Automatic Purity}
\label{sec:face-gap}

\begin{definition}
The \textbf{face dimension spectrum} of $\mathcal{A}$ is
$\Spec(\mathcal{A}) = \{\dim(F_{S(\mathcal{A})}(\varphi)) :
\varphi \in S(\mathcal{A})\}$.
The \textbf{face dimension gap} is
$G(\mathcal{A}) = \{d \in \N_{>0} : d \notin \Spec(\mathcal{A})\}$.
\end{definition}

\begin{theorem}[Automatic Purity]
\label{thm:automatic-purity}
If $G(\mathcal{A}) \supseteq \{1, 2, \ldots, \ell - 1\}$ (i.e., the
smallest positive face dimension is $\ge \ell$), then for any $C$ defined
by $m \le \ell - 1$ affine constraints, every extreme point of $C$ is a
\emph{pure state}.
\end{theorem}

\begin{proof}
Let $\varphi$ be extreme in $C$. By \cref{thm:main},
$\dim(F_{S(\mathcal{A})}(\varphi)) \le m \le \ell - 1$.
By the gap hypothesis,
$\dim(F_{S(\mathcal{A})}(\varphi)) \notin \{1, \ldots, \ell - 1\}$.
Since $\dim \ge 0$, the only possibility is $\dim = 0$, i.e.,
$\varphi$ is pure.
\end{proof}

\begin{example}
For $\BH$: $\Spec = \{0, 3, 8, 15, \ldots\}$
(face dimensions $r^2 - 1$ for $r = 1, 2, 3, \ldots$), so
$G \supseteq \{1, 2\}$ and $\ell = 3$. Under $m \le 2$ energy
constraints, all extreme points are pure states.
This is the content of Weis--Shirokov Corollary~7.
\end{example}

The face dimension spectrum depends on the algebra type:

\begin{center}
\begin{tabular}{lll}
\toprule
\textbf{Algebra} & $\Spec$ & \textbf{Gap} \\
\midrule
$C(X)$, $\abs{X} = \infty$ & $\{0, 1, 2, \ldots\}$ & $\emptyset$ \\
$M_n(\C)$, $n \ge 2$ & $\{0, 3, 8, 15, \ldots\}$ & $\supseteq \{1, 2\}$ \\
$M_2(\C) \oplus M_2(\C)$ & $\{0, 1, 3, 4, 7\}$ & $\{2, 5, 6\}$ \\
$\BH$, $\dim\Hilbert = \infty$ & $\{0, 3, 8, \ldots\} \cup \{\infty\}$
& $\supseteq \{1, 2\}$ \\
\bottomrule
\end{tabular}
\end{center}


%% ============================================================
%% 6. APPLICATIONS
%% ============================================================
\section{Applications}
\label{sec:applications}

\subsection{NPA Hierarchy}
\label{sec:npa}

At finite level $k$ of the NPA hierarchy, the feasible set $F_k$ is
an SDP spectrahedron: PSD moment matrices $\Gamma \in M_{N_k}(\C)$
with $m_k$ affine constraints. The main theorem (applied with
$\mathcal{A} = M_{N_k}(\C)$) gives
$\rank(\Gamma)^2 \le m_k + 1$ for extreme points.

Beyond the rank bound, the theorem provides structural information:
the GNS decomposition $\sum n_j^2 \le m_k + 1$ reveals the multiplicity
structure of extremal solutions, which classical Barvinok--Pataki
does not capture.

\begin{caveat*}
Interior-point SDP solvers converge to the \emph{analytic centre}
of the optimal face, which has \emph{maximal} rank and is typically
\emph{not} an extreme point. Post-processing (facial reduction or
Burer--Monteiro methods) is needed to obtain low-rank extreme points.
\end{caveat*}

\subsection{Quantum Marginal Problem}
\label{sec:quantum-marginals}

Let $\Hilbert = \bigotimes_{i \in I} \Hilbert_i$ be a finite-dimensional
multipartite Hilbert space with non-overlapping subsystems
$S_1, \ldots, S_k \subseteq I$. The compatible state set is
$C = \{\rho \in S(\BH) : \tr_{I \setminus S_j}(\rho) = \rho_j\}$.

Each marginal constraint decomposes into $D_j^2 - 1$ independent scalar
affine constraints (where $D_j = \dim(\Hilbert_{S_j})$), since the trace
component is automatically satisfied by the normalisation $\varphi(1) = 1$.
For non-overlapping subsystems, $m = \sum_{j=1}^k (D_j^2 - 1)$.

For qubit marginals ($D_j = 2$): $m = 3k$, and the bound is
$\sum n_i^2 \le 3k + 1$.


%% ============================================================
%% 7. ERRORS FOUND DURING VERIFICATION
%% ============================================================
\section{Errors Found During Adversarial Verification}
\label{sec:errors}

The adversarial verification process discovered four genuine mathematical
errors in the original proof draft:

\begin{enumerate}
  \item \textbf{Face dimension spectrum} (Node 1.9.2):
        The spectrum of $M_2(\C) \oplus M_2(\C)$ is $\{0, 1, 3, 4, 7\}$,
        \emph{not} $\{0, 3, 7\}$ as originally claimed. States supported
        on two blocks with rank~1 in each create face dimension~1.
        The gap $\{1, 2\}$ holds only for \emph{simple} algebras ($N = 1$).

  \item \textbf{``Trivial involution''} (Node 1.6.6):
        The real PSD cone uses the \emph{transpose} involution $a^* = a^\top$,
        which is a genuine anti-automorphism, not the identity.
        The notation $\mathcal{A} = M_n(\R)^{\sa}$ was incorrect:
        the algebra is $M_n(\R)$ with transpose involution; the
        self-adjoint part $M_n(\R)^{\sa} = S_n$ is the order unit space,
        not the algebra.

  \item \textbf{KMV scope mismatch} (Node 1.10.3):
        The Klep--Magron--Vol\v{c}i\v{c} state polynomial constraints
        are \emph{nonlinear} in the state $\varphi$, but the main theorem
        handles only \emph{affine} constraints. The application was
        restricted to the linear sub-case.

  \item \textbf{Marginal constraint count} (Node 1.10.4):
        Each marginal contributes $\dim(\Hilbert_{S_j})^2 - 1$
        constraints, not $\dim(\Hilbert_{S_j})^2$. The trace constraint
        is automatically satisfied by normalisation.
        For qubits: $m = 3k$, not $4k$.
\end{enumerate}


%% ============================================================
%% 8. VERIFICATION METHODOLOGY
%% ============================================================
\section{Verification Methodology}
\label{sec:methodology}

The proof was constructed and verified using the \texttt{af}
(Adversarial Proof Framework) CLI tool, which implements a
prover--verifier protocol with strict role isolation.

\begin{itemize}
  \item \textbf{Proof tree:} 64 nodes across 3 levels of depth,
        with 32 definitions and 14 external references.
  \item \textbf{Verification:} 53 leaf nodes were independently
        verified by opus-level adversarial agents.
        15 nodes were accepted outright; 38 were challenged with
        148 total challenges.
  \item \textbf{Challenge resolution:} All 148 challenges have been
        resolved. Of these, 4 identified genuine mathematical errors
        (corrected via \texttt{af amend}), 6 identified critical
        logical gaps (filled with additional proof steps), and the
        remainder were structural metadata issues (missing dependency
        declarations, incorrect inference types).
  \item \textbf{Re-verification pending:} The 49 amended nodes
        require a second verifier pass before the proof can be
        considered fully validated.
\end{itemize}

\begin{center}
\begin{tabular}{lr}
\toprule
\textbf{Metric} & \textbf{Value} \\
\midrule
Total nodes & 64 \\
Validated & 15 (23\%) \\
Pending & 49 (77\%) \\
Challenges raised & 148 \\
Challenges resolved & 148 (100\%) \\
Genuine errors found & 4 \\
Critical gaps found & 6 \\
Quality score & 69.4/100 \\
\bottomrule
\end{tabular}
\end{center}


%% ============================================================
%% REFERENCES
%% ============================================================
\begin{thebibliography}{99}

\bibitem{alfsen2001}
E.\,M.\ Alfsen and F.\,W.\ Shultz,
\emph{State Spaces of Operator Algebras: Basic Theory, Orientations,
and $C^*$-Products},
Birkh\"auser, 2001.

\bibitem{alfsen2003}
E.\,M.\ Alfsen and F.\,W.\ Shultz,
\emph{Geometry of State Spaces of Operator Algebras},
Birkh\"auser, 2003.

\bibitem{barvinok2001}
A.\ Barvinok,
A remark on the rank of positive semidefinite matrices subject to
affine constraints,
\emph{Discrete Comput.\ Geom.}\ \textbf{25} (2001), 23--31.

\bibitem{klep2023}
I.\ Klep, V.\ Magron, and J.\ Vol\v{c}i\v{c},
State polynomials: positivity, optimization and nonlinear Bell
inequalities,
\emph{Math.\ Program.}\ (2023).

\bibitem{navascues2008}
M.\ Navascu\'es, S.\ Pironio, and A.\ Ac\'in,
A convergent hierarchy of semidefinite programs characterizing the
set of quantum correlations,
\emph{New J.\ Phys.}\ \textbf{10} (2008), 073013.

\bibitem{pataki1998}
G.\ Pataki,
On the rank of extreme matrices in semidefinite programs and the
multiplicity of optimal eigenvalues,
\emph{Math.\ Oper.\ Res.}\ \textbf{23} (1998), 339--358.

\bibitem{pataki2000}
G.\ Pataki,
Geometry of semidefinite programming,
in \emph{Handbook of Semidefinite Programming} (H.\ Wolkowicz et al.,
eds.), Kluwer, 2000, pp.\ 29--65.

\bibitem{weisshirokov2021}
S.\ Weis and M.\,E.\ Shirokov,
The face generated by a point, generalized affine constraints, and
quantum theory,
\emph{J.\ Convex Anal.}\ \textbf{28}(3) (2021), 847--870.

\end{thebibliography}


%% ============================================================
%% APPENDIX: PROOF TREE
%% ============================================================
\newpage
\appendix

\section{Proof Tree}
\label{app:proof-tree}

The complete proof tree, exported from the \texttt{af} framework, is
shown below. Each node is annotated with its epistemic state
(\textsc{validated} or \textsc{pending}) and the number of challenges
raised and resolved.

\begin{tcolorbox}[colback=blue!5!white, colframe=blue!50!black,
  title={\textbf{Proof Tree Legend}}]
\begin{itemize}[nosep]
  \item \textbf{Node ID}: hierarchical numbering (e.g., 1.3.2)
  \item \textbf{Type}: \texttt{claim} (assertion requiring proof)
  \item \textbf{Status}: \textsc{validated} (accepted by verifier)
        or \textsc{pending} (awaiting verification)
  \item Challenges shown as {[raised/resolved]}
\end{itemize}
\end{tcolorbox}

\subsection{Root and Level-1 Nodes}

\begin{longtable}{llll}
\toprule
\textbf{Node} & \textbf{Title} & \textbf{Status} & \textbf{Challenges} \\
\midrule
\endfirsthead
\toprule
\textbf{Node} & \textbf{Title} & \textbf{Status} & \textbf{Challenges} \\
\midrule
\endhead
\bottomrule
\endlastfoot
1     & Main Theorem (BP for OUS) & pending & --- \\
1.1   & Face-Commutant (Alfsen--Shultz) & pending & --- \\
1.2   & Face-Dimension Bound (Weis--Shirokov) & pending & --- \\
1.3   & Main Dimension Bound & pending & --- \\
1.4   & Commutant Decomposition & pending & --- \\
1.5   & Recovery: Abelian $\to$ Carath\'eodory & pending & --- \\
1.6   & Recovery: Matrix $\to$ Classical BP & pending & --- \\
1.7   & Recovery: Block-Diagonal $\to$ Mixed & pending & --- \\
1.8   & Recovery: $\BH$ $\to$ Weis--Shirokov & pending & --- \\
1.9   & Face Dimension Gap \& Purity & pending & --- \\
1.10  & Applications (NPA, KMV, Marginals) & pending & --- \\
1.11  & Consistency Check & pending & --- \\
\end{longtable}

\subsection{Level-2 Nodes: Core Proof}

\begin{longtable}{lp{6cm}ll}
\toprule
\textbf{Node} & \textbf{Statement (abridged)} & \textbf{Status} & \textbf{Ch.} \\
\midrule
\endfirsthead
\toprule
\textbf{Node} & \textbf{Statement (abridged)} & \textbf{Status} & \textbf{Ch.} \\
\midrule
\endhead
\bottomrule
\endlastfoot
1.1.1 & Well-definedness of $\Phi$ & pending & 1/1 \\
1.1.2 & $\Phi$ is affine & pending & 2/2 \\
1.1.3 & Image lies in face & pending & 2/2 \\
1.1.4 & Injectivity & validated & 1/1 \\
1.1.5 & Surjectivity (Radon--Nikodym) & pending & 4/4 \\
1.1.6 & Dimension consequence & pending & 1/1 \\
\midrule
1.2.1 & Single-constraint (WS Thm 2) & pending & 1/1 \\
1.2.2 & Multi-constraint induction (WS Cor 5) & pending & 3/3 \\
1.2.3 & Extremal specialisation & pending & 4/4 \\
1.2.4 & Generalised vs ordinary affine maps & pending & 2/2 \\
\midrule
1.3.1 & Hypothesis verification (corrected scope) & pending & 6/6 \\
1.3.2 & Step 1: geometric-algebraic bridge & pending & 3/3 \\
1.3.3 & Step 2: extremality $\Rightarrow$ face bound & pending & 5/5 \\
1.3.4 & Step 3: conclusion & pending & 2/2 \\
1.3.5 & Extension to inequality/mixed constraints & pending & 2/2 \\
\midrule
1.4.1 & Finite-dimensionality of commutant & pending & 1/1 \\
1.4.2 & Artin--Wedderburn decomposition & pending & 3/3 \\
1.4.3 & Correspondence with isotypic components & pending & 5/5 \\
1.4.4 & Dimension calculation $\sum n_j^2$ & pending & 3/3 \\
1.4.5 & Contrapositive for infinite-dim case & pending & 3/3 \\
\end{longtable}

\subsection{Level-2 Nodes: Recovery and Applications}

\begin{longtable}{lp{6cm}ll}
\toprule
\textbf{Node} & \textbf{Statement (abridged)} & \textbf{Status} & \textbf{Ch.} \\
\midrule
\endfirsthead
\toprule
\textbf{Node} & \textbf{Statement (abridged)} & \textbf{Status} & \textbf{Ch.} \\
\midrule
\endhead
\bottomrule
\endlastfoot
1.5.1 & Abelian setup: $C(X)$, states = measures & validated & 2/2 \\
1.5.2 & All irreps 1-dimensional & pending & 2/2 \\
1.5.3 & GNS decomposition for measures & pending & 4/4 \\
1.5.4 & Applying the bound: $k \le m+1$ & pending & 3/3 \\
\midrule
1.6.1 & Matrix algebra setup: $M_n(\C)$ & validated & 1/1 \\
1.6.2 & Unique irreducible representation & pending & 3/3 \\
1.6.3 & GNS of rank-$r$ density matrix & pending & 2/2 \\
1.6.4 & Applying the bound: $r^2 \le m+1$ & pending & 3/3 \\
1.6.5 & $+1$ reconciliation with classical BP & validated & 2/2 \\
1.6.6 & Real vs complex subtlety (corrected) & pending & 3/3 \\
\midrule
1.7.1 & Block-diagonal setup & validated & 1/1 \\
1.7.2 & GNS decomposition & pending & 3/3 \\
1.7.3 & Applying the bound: $\sum r_i^2 \le m+1$ & pending & 2/2 \\
1.7.4 & Novelty of the mixed bound & pending & 4/4 \\
\midrule
1.8.1 & Infinite-dimensional setup & pending & 3/3 \\
1.8.2 & GNS for $\BH$ & pending & 3/3 \\
1.8.3 & Energy constraints (corrected) & pending & 4/4 \\
1.8.4 & $m \le 2$: automatic purity & pending & 3/3 \\
1.8.5 & Extension beyond Type I & validated & 6/6 \\
\midrule
1.9.1 & Definition of spectrum and gap & pending & 4/4 \\
1.9.2 & Computing the gap (corrected) & pending & 2/2 \\
1.9.3 & Automatic purity theorem & pending & 3/3 \\
1.9.4 & Application to $\BH$ & pending & 4/4 \\
1.9.5 & Quantitative gap formula & pending & 4/4 \\
\midrule
1.10.1 & NPA hierarchy (corrected) & pending & 5/5 \\
1.10.2 & Structural consequence (corrected) & pending & 3/3 \\
1.10.3 & KMV state polynomials (scope corrected) & pending & 6/6 \\
1.10.4 & Quantum marginals (count corrected) & pending & 7/7 \\
\midrule
1.11.1 & Bloch ball consistency & validated & 0/0 \\
1.11.2 & Maximally mixed state & validated & 2/2 \\
1.11.3 & Why 3 constraints give non-pure & validated & 0/0 \\
1.11.4 & Consistency with Im 2025 & validated & 0/0 \\
\end{longtable}

\subsection{Full Export}

The raw \texttt{af export --format latex} output is available in the
companion file \texttt{proof-tree-export.tex} (78\,KB).

\end{document}
